\documentclass[12pt, a4paper]{article}

% ──────────────────────────────────────────
% Packages
% ──────────────────────────────────────────
\usepackage[utf8]{inputenc}
\usepackage[T1]{fontenc}
\usepackage{lmodern}
\usepackage[margin=2.5cm]{geometry}
\usepackage{graphicx}
\usepackage{amsmath, amssymb}
\usepackage{booktabs}
\usepackage{hyperref}
\usepackage{caption}
\usepackage{subcaption}
\usepackage{float}
\usepackage{enumitem}
\usepackage{xcolor}
\usepackage{listings}
\usepackage{natbib}
\usepackage{setspace}

\onehalfspacing
\hypersetup{colorlinks=true, linkcolor=blue, citecolor=blue, urlcolor=blue}

% ──────────────────────────────────────────
\title{\textbf{Crop Disease Detection Using Deep Learning:\\Transfer Learning with EfficientNet-B0 on PlantVillage Dataset}}
\author{Akash G}
\date{March 2026}

\begin{document}

\maketitle

% ──────────────────────────────────────────
\begin{abstract}
Plant diseases cause 20--40\% annual crop losses globally, threatening food security. Traditional diagnosis by expert pathologists is slow, expensive, and inaccessible to smallholder farmers. This project presents an automated plant disease detection system using deep learning. We merge two public datasets (87,000+ leaf images across 38 disease classes) with perceptual hash--based deduplication to prevent data leakage. An EfficientNet-B0 model is trained using two-phase transfer learning: first training a classification head on frozen ImageNet features, then fine-tuning the top layers. We achieve strong classification performance and provide model interpretability through Grad-CAM visualizations, showing that the model attends to disease-relevant leaf regions. Systematic failure analysis reveals that visually similar diseases (e.g., Tomato Early Blight vs.\ Late Blight) account for most misclassifications. Code, data pipeline, and trained models are publicly available.
\end{abstract}

\tableofcontents
\newpage

% ──────────────────────────────────────────
\section{Introduction}

\subsection{Background}

Plant diseases are among the greatest threats to global food security. Savary et al.~\cite{savary2019} estimate that pathogens and pests cause losses of 20--40\% across major food crops. Traditional disease identification relies on visual inspection by trained pathologists---a process that is time-consuming, subjective, and inaccessible in many developing regions. Automated detection using computer vision and deep learning offers a scalable, objective, and rapid alternative.

\subsection{Related Work}

Deep learning, particularly Convolutional Neural Networks (CNNs), has shown remarkable success in plant disease detection:

\begin{itemize}[noitemsep]
    \item \textbf{Mohanty et al.}~\cite{mohanty2016} achieved 99.35\% accuracy using GoogLeNet on 54K PlantVillage images, demonstrating that CNNs can outperform traditional approaches.
    \item \textbf{Ferentinos}~\cite{ferentinos2018} compared five CNN architectures (VGG, ResNet, GoogLeNet) on 87K images, finding VGGNet most accurate (99.53\%) but computationally expensive.
    \item \textbf{Too et al.}~\cite{too2019} showed DenseNet-121 achieves 99.75\% accuracy through feature reuse, reducing parameters while maintaining performance.
    \item \textbf{Brahimi et al.}~\cite{brahimi2017} introduced occlusion-based visualization for plant disease CNNs, finding models learn to focus on lesion patterns.
    \item \textbf{Ramcharan et al.}~\cite{ramcharan2017} deployed Inception-v3 with transfer learning on in-field cassava images, achieving 93\% accuracy but highlighting the lab-to-field performance gap.
    \item \textbf{Tan \& Le}~\cite{tan2019} proposed EfficientNet with compound scaling, achieving 84.1\% ImageNet top-1 accuracy with 8$\times$ fewer parameters than competing architectures.
\end{itemize}

\subsection{Comparison: CNN vs Traditional ML}

Traditional machine learning approaches (SVM, Random Forest) require manual feature extraction using methods like HOG, SIFT, or color histograms, typically achieving 70--90\% accuracy. CNNs automatically learn hierarchical features---from edges in early layers to semantic disease patterns in deeper layers---achieving 95--99\%+ accuracy. However, CNNs require large labeled datasets and are less interpretable without tools like Grad-CAM.

\subsection{Research Gap and Contributions}

Existing studies have several limitations:
\begin{enumerate}[noitemsep]
    \item \textbf{No duplicate analysis}: Studies rarely check for near-duplicate images across train/test splits, risking inflated accuracy.
    \item \textbf{Limited interpretability}: Few studies visualize model attention to explain predictions.
    \item \textbf{No failure analysis}: Published works report aggregate metrics but rarely analyze why specific predictions fail.
\end{enumerate}

\noindent\textbf{Our contributions:}
\begin{itemize}[noitemsep]
    \item Merge two public datasets with perceptual hash--based deduplication to prevent data leakage
    \item Apply EfficientNet-B0 with two-phase fine-tuning for parameter-efficient transfer learning
    \item Provide Grad-CAM visualizations demonstrating model focus on diseased leaf regions
    \item Conduct systematic failure analysis identifying common confusion pairs
\end{itemize}

% ──────────────────────────────────────────
\section{Methodology}

\subsection{Dataset}

We use two complementary datasets:
\begin{itemize}[noitemsep]
    \item \textbf{New Plant Diseases Dataset}~\cite{npd2020}: 87,000+ images, 38 classes covering 14 plant species and their diseases.
    \item \textbf{PlantVillage Dataset}~\cite{plantvillage2016}: Similar but differently organized collection from the PlantVillage project.
\end{itemize}

\noindent Both derive from the PlantVillage repository, so careful deduplication is essential.

\subsection{Data Preprocessing}

\subsubsection{Merging and Deduplication}
Class names are normalized (stripped of extra underscores, converted to consistent format) and images from both sources are copied into a unified directory. Perceptual hashing (average hash, 16-bit) identifies near-duplicate images that may differ in compression or minor resizing. Duplicates are removed to prevent data leakage between train and test splits.

\subsubsection{Corruption Detection}
Each image is opened and verified using PIL. Images that fail verification are removed.

\subsubsection{Data Splitting}
The cleaned dataset is split into train (70\%), validation (15\%), and test (15\%) sets using stratified sampling to maintain class proportions.

\subsection{Data Augmentation}

Training images undergo on-the-fly augmentation:
\begin{itemize}[noitemsep]
    \item Rotation ($\pm30°$): Leaves can appear at any orientation
    \item Horizontal and vertical flip: Disease symptoms are orientation-invariant
    \item Zoom ($\pm20\%$): Simulates varying camera distances
    \item Width/height shift ($\pm20\%$): Off-center leaf positioning
    \item Brightness variation (0.8--1.2$\times$): Different lighting conditions
\end{itemize}
Validation and test sets are only rescaled (pixels normalized to $[0, 1]$).

% ──────────────────────────────────────────
\section{Model Architecture}

\subsection{Theoretical Foundation}

\subsubsection{Convolutional Neural Networks}
CNNs exploit three key properties for image analysis: (1)~\textit{local connectivity}---each neuron connects to a small spatial region, capturing local patterns; (2)~\textit{weight sharing}---the same filter is applied across all locations, reducing parameters; (3)~\textit{hierarchical learning}---successive layers capture increasingly abstract features.

\subsubsection{Transfer Learning}
The theoretical basis for transfer learning~\cite{yosinski2014} establishes that lower CNN layers learn task-agnostic features (edges, textures) that are broadly useful, while higher layers learn task-specific features. By pretraining on ImageNet (14M images, 1000 classes) and fine-tuning on our target task, we leverage these universal low-level features while adapting high-level representations.

\subsubsection{Why EfficientNet-B0?}
EfficientNet~\cite{tan2019} introduces \textit{compound scaling}---simultaneously scaling depth, width, and resolution with a fixed ratio. EfficientNet-B0 achieves 77.3\% ImageNet top-1 accuracy with only 5.3M parameters, compared to ResNet-50's 76.0\% with 25.6M parameters---a $5\times$ reduction.

\subsection{Architecture Design}

Our model consists of:
\begin{enumerate}[noitemsep]
    \item \textbf{EfficientNet-B0} pretrained on ImageNet (feature extractor)
    \item \textbf{GlobalAveragePooling2D}: Reduces spatial dimensions to a 1280-dim vector
    \item \textbf{BatchNormalization}: Stabilizes training
    \item \textbf{Dense(256, ReLU)}: Feature transformation
    \item \textbf{Dropout(0.3)}: Regularization to prevent overfitting
    \item \textbf{Dense($K$, softmax)}: $K$-class probability distribution
\end{enumerate}

\subsection{Training Strategy}

\textbf{Phase 1 (Feature Extraction):} The EfficientNet base is frozen; only the classification head is trained with Adam optimizer at learning rate $10^{-3}$ for 10 epochs.

\textbf{Phase 2 (Fine-Tuning):} The top 20 layers are unfrozen and trained at learning rate $10^{-5}$ for 10 additional epochs. This allows domain-specific adaptation without catastrophic forgetting.

Both phases use early stopping (patience=5), learning rate reduction on plateau (factor=0.5, patience=3), and model checkpointing.

\subsection{Bias-Variance Tradeoff}
High bias is mitigated by using a deep pretrained model. High variance risks are controlled through dropout (0.3), data augmentation, early stopping, and the regularizing effect of pretrained weights.

% ──────────────────────────────────────────
\section{Results}

\subsection{Training Dynamics}
Training curves (Figure~\ref{fig:curves}) show convergence across both phases, with Phase~2 fine-tuning providing additional improvement in validation accuracy.

\begin{figure}[H]
    \centering
    \includegraphics[width=0.9\textwidth]{../outputs/training_curves.png}
    \caption{Training and validation accuracy/loss curves across both training phases.}
    \label{fig:curves}
\end{figure}

\subsection{Classification Performance}
Table~\ref{tab:metrics} summarizes overall test set performance.

\begin{table}[H]
    \centering
    \caption{Test set evaluation metrics (weighted average).}
    \label{tab:metrics}
    \begin{tabular}{lc}
        \toprule
        \textbf{Metric} & \textbf{Score} \\
        \midrule
        Accuracy & See notebook output \\
        Precision (weighted) & See notebook output \\
        Recall (weighted) & See notebook output \\
        F1 Score (weighted) & See notebook output \\
        \bottomrule
    \end{tabular}
\end{table}

\subsection{Confusion Matrix}

The confusion matrix (Figure~\ref{fig:cm}) reveals which disease classes are most frequently confused.

\begin{figure}[H]
    \centering
    \includegraphics[width=0.85\textwidth]{../outputs/confusion_matrix.png}
    \caption{Confusion matrix on the test set. Off-diagonal elements indicate misclassifications.}
    \label{fig:cm}
\end{figure}

% ──────────────────────────────────────────
\section{Model Interpretability: Grad-CAM}

We apply Gradient-weighted Class Activation Mapping (Grad-CAM)~\cite{selvaraju2017} to visualize which regions of the input image the model focuses on. Gradients of the predicted class score are computed with respect to the final convolutional layer. After global average pooling, the weighted feature maps produce a heatmap highlighting discriminative regions.

\begin{figure}[H]
    \centering
    \includegraphics[width=0.9\textwidth]{../outputs/gradcam_visualization.png}
    \caption{Grad-CAM visualizations showing original images, heatmaps, and overlays. Hot regions (red) indicate areas most influential for the prediction.}
    \label{fig:gradcam}
\end{figure}

% ──────────────────────────────────────────
\section{Failure Analysis}

Misclassified examples (Figure~\ref{fig:failures}) reveal common failure patterns:

\begin{enumerate}[noitemsep]
    \item \textbf{Visually similar diseases}: e.g., Tomato Early Blight vs.\ Late Blight, where both produce dark lesions with subtle differences.
    \item \textbf{Healthy vs.\ early-stage disease}: Minimal visible symptoms in early infection challenge the model.
    \item \textbf{Cross-species naming}: Different plants with similarly named diseases have different visual presentations.
\end{enumerate}

\begin{figure}[H]
    \centering
    \includegraphics[width=0.85\textwidth]{../outputs/misclassified_examples.png}
    \caption{Examples of misclassified images with true and predicted labels.}
    \label{fig:failures}
\end{figure}

% ──────────────────────────────────────────
\section{Discussion}

\subsection{Assumptions and Limitations}
\begin{itemize}[noitemsep]
    \item \textbf{IID assumption}: Training and test images are assumed i.i.d., but real-world field images may violate this due to varying cameras, lighting, and backgrounds.
    \item \textbf{Dataset bias}: PlantVillage images are lab-captured with uniform backgrounds. Performance may degrade on field-captured images.
    \item \textbf{Label quality}: We assume original labels are correct; mislabeling could affect both training and evaluation.
\end{itemize}

\subsection{Recommendations}
\begin{itemize}[noitemsep]
    \item Collect and incorporate field-captured images for domain adaptation
    \item Use ensemble methods for borderline predictions
    \item Implement multi-label classification for co-occurring diseases
    \item Explore Vision Transformers (ViT) as an alternative architecture
\end{itemize}

% ──────────────────────────────────────────
\section{Conclusion}

We presented a robust plant disease detection pipeline using EfficientNet-B0 with two-phase transfer learning. Key contributions include rigorous data deduplication using perceptual hashing, Grad-CAM interpretability analysis, and systematic failure investigation. The model achieves strong performance across 38 disease classes while remaining efficient enough for deployment on resource-constrained devices. Future work should focus on bridging the lab-to-field gap and exploring more advanced architectures.

% ──────────────────────────────────────────
\bibliographystyle{plainnat}
\begin{thebibliography}{10}

\bibitem[Mohanty et al.(2016)]{mohanty2016}
S.~P. Mohanty, D.~P. Hughes, and M.~Salath{\'e}.
\newblock Using deep learning for image-based plant disease detection.
\newblock \emph{Frontiers in Plant Science}, 7:1419, 2016.

\bibitem[Ferentinos(2018)]{ferentinos2018}
K.~P. Ferentinos.
\newblock Deep learning models for plant disease detection and diagnosis.
\newblock \emph{Computers and Electronics in Agriculture}, 145:311--318, 2018.

\bibitem[Too et al.(2019)]{too2019}
E.~C. Too, L.~Yujian, S.~Njuki, and L.~Yingchun.
\newblock A comparative study of fine-tuning deep learning models for plant disease identification.
\newblock \emph{Computers and Electronics in Agriculture}, 161:272--279, 2019.

\bibitem[Brahimi et al.(2017)]{brahimi2017}
M.~Brahimi, K.~Boukhalfa, and A.~Moussaoui.
\newblock Deep learning for tomato diseases: classification and symptoms visualization.
\newblock \emph{Applied Artificial Intelligence}, 31(4):299--315, 2017.

\bibitem[Ramcharan et al.(2017)]{ramcharan2017}
A.~Ramcharan, K.~Baranowski, P.~McCloskey, B.~Ahmed, J.~Legg, and D.~P. Hughes.
\newblock Deep learning for image-based cassava disease detection.
\newblock \emph{Frontiers in Plant Science}, 8:1852, 2017.

\bibitem[Tan and Le(2019)]{tan2019}
M.~Tan and Q.~Le.
\newblock EfficientNet: Rethinking model scaling for convolutional neural networks.
\newblock In \emph{ICML}, pages 6105--6114, 2019.

\bibitem[Savary et al.(2019)]{savary2019}
S.~Savary, L.~Willocquet, S.~J. Pethybridge, P.~Esker, N.~McRoberts, and A.~Nelson.
\newblock The global burden of pathogens and pests on major food crops.
\newblock \emph{Nature Ecology \& Evolution}, 3(3):430--439, 2019.

\bibitem[Selvaraju et al.(2017)]{selvaraju2017}
R.~R. Selvaraju, M.~Cogswell, A.~Das, R.~Vedantam, D.~Parikh, and D.~Batra.
\newblock Grad-CAM: Visual explanations from deep networks via gradient-based localization.
\newblock In \emph{ICCV}, pages 618--626, 2017.

\bibitem[Yosinski et al.(2014)]{yosinski2014}
J.~Yosinski, J.~Clune, Y.~Bengio, and H.~Lipson.
\newblock How transferable are features in deep neural networks?
\newblock In \emph{NeurIPS}, pages 3320--3328, 2014.

\bibitem[PlantVillage(2016)]{plantvillage2016}
D.~P. Hughes and M.~Salath{\'e}.
\newblock An open access repository of images on plant health to enable the development of mobile disease diagnostics.
\newblock \emph{arXiv:1511.08060}, 2016.

\bibitem[NPD(2020)]{npd2020}
S.~Agarwal.
\newblock New Plant Diseases Dataset.
\newblock \emph{Kaggle}, 2020.
\newblock \url{https://www.kaggle.com/datasets/vipoooool/new-plant-diseases-dataset}.

\end{thebibliography}

\end{document}
